\documentclass[11pt,a4paper]{article}

%========================
% PACKAGES
%========================
\usepackage[utf8]{inputenc}
\usepackage[T1]{fontenc}
\usepackage[french]{babel}
\usepackage[a4paper,margin=2cm]{geometry}
\usepackage{amsmath,amssymb,amsfonts,amsthm}
\usepackage{mathtools}
\usepackage{lmodern}
\usepackage{microtype}
\usepackage{xcolor}
\usepackage{enumitem}
\usepackage{fancyhdr}
\usepackage{lastpage}
\usepackage{tikz}
\usepackage{hyperref}

%========================
% STYLE
%========================
\definecolor{bleu}{RGB}{20,70,140}
\definecolor{rouge}{RGB}{160,30,30}
\definecolor{vert}{RGB}{20,120,60}
\definecolor{gris}{RGB}{245,245,245}

\hypersetup{
    colorlinks=true,
    linkcolor=bleu,
    urlcolor=bleu,
    citecolor=bleu
}

\setlength{\headheight}{14pt}
\pagestyle{fancy}
\fancyhf{}
\fancyhead[L]{Pré-requis}
\fancyhead[C]{Avant un cours de topologie générale}
\fancyhead[R]{Page \thepage/\pageref{LastPage}}

%========================
% ENVIRONNEMENTS
%========================
\theoremstyle{definition}
\newtheorem{definition}{Définition}[section]
\newtheorem{exemple}[definition]{Exemple}
\newtheorem{remarque}[definition]{Remarque}
\newtheorem{methode}[definition]{Méthode}

\theoremstyle{plain}
\newtheorem{proposition}[definition]{Proposition}
\newtheorem{theoreme}[definition]{Théorème}
\newtheorem{lemme}[definition]{Lemme}
\newtheorem{corollaire}[definition]{Corollaire}

\newcommand{\R}{\mathbb{R}}
\newcommand{\N}{\mathbb{N}}
\newcommand{\Q}{\mathbb{Q}}
\newcommand{\C}{\mathbb{C}}
\newcommand{\K}{\mathbb{K}}
\newcommand{\eps}{\varepsilon}
\newcommand{\id}{\mathrm{Id}}
\newcommand{\card}{\mathrm{Card}}
\newcommand{\adh}{\overline}
\newcommand{\interieur}[1]{\mathring{#1}}
\newcommand{\norme}[1]{\left\lVert #1 \right\rVert}
\newcommand{\abs}[1]{\left\lvert #1 \right\rvert}
\newcommand{\dist}{\mathrm{d}}

%========================
% DOCUMENT
%========================
\begin{document}

\begin{center}
    {\Huge \textbf{Pré-requis mathématiques}}\\[0.2cm]
    {\Large \textbf{Avant de comprendre un cours de topologie générale}}\\[0.3cm]
    {\large Niveau : bac+3}\\[0.3cm]
    \rule{0.85\textwidth}{0.5pt}
\end{center}

\vspace{0.5cm}

\tableofcontents

\newpage

%================================================
\section*{Objectif du document}
\addcontentsline{toc}{section}{Objectif du document}

Ce document rassemble les notions qu'il faut maîtriser avant d'aborder sérieusement un cours de topologie générale.

Le but n'est pas encore de refaire tout le cours de topologie, mais de préparer les outils indispensables :
\begin{itemize}
    \item logique et quantificateurs ;
    \item ensembles et applications ;
    \item relations d'équivalence et quotients ;
    \item suites et limites ;
    \item espaces vectoriels normés ;
    \item normes, distances, boules ;
    \item continuité ;
    \item compacité dans $\R^n$ ;
    \item connexité intuitive ;
    \item suites de Cauchy et complétude.
\end{itemize}

Le cours de topologie générale abstrait toutes ces notions. Par exemple, les ouverts ne seront plus forcément des intervalles ou des boules : ils seront définis par une famille $\tau$ de parties d'un ensemble $E$ vérifiant certaines propriétés.

\newpage

%================================================
\section{Logique, implications et quantificateurs}

\subsection{Propositions mathématiques}

\begin{definition}
Une \textbf{proposition} est une phrase mathématique ayant une valeur de vérité : vraie ou fausse.
\end{definition}

\begin{exemple}
Les phrases suivantes sont des propositions :
\[
2+2=4,
\qquad
\forall x\in\R,\ x^2\geq 0,
\qquad
\exists x\in\R,\ x^2=-1.
\]
La première et la deuxième sont vraies, la troisième est fausse.
\end{exemple}

\subsection{Implication}

\begin{definition}
L'implication
\[
P \Rightarrow Q
\]
signifie : \og si $P$ est vraie, alors $Q$ est vraie \fg.
\end{definition}

\begin{remarque}
Une implication n'affirme pas que $P$ est vraie. Elle affirme seulement que, dans le cas où $P$ est vraie, alors $Q$ l'est aussi.
\end{remarque}

\begin{exemple}
\[
x>2 \Rightarrow x>0.
\]
Cette implication est vraie. Elle ne dit pas que tous les réels sont strictement supérieurs à $2$.
\end{exemple}

\subsection{Réciproque, contraposée, négation}

\begin{definition}
À partir de l'implication $P\Rightarrow Q$, on définit :
\begin{itemize}
    \item la \textbf{réciproque} : $Q\Rightarrow P$ ;
    \item la \textbf{contraposée} : $\neg Q\Rightarrow \neg P$ ;
    \item la \textbf{négation} : $P$ est vraie et $Q$ est fausse.
\end{itemize}
\end{definition}

\begin{proposition}
Une implication est équivalente à sa contraposée :
\[
(P\Rightarrow Q)\quad \Longleftrightarrow \quad (\neg Q\Rightarrow \neg P).
\]
\end{proposition}

\begin{exemple}
L'implication
\[
x>2 \Rightarrow x>0
\]
a pour contraposée
\[
x\leq 0 \Rightarrow x\leq 2.
\]
Elle est vraie.

Sa réciproque est
\[
x>0 \Rightarrow x>2,
\]
qui est fausse.
\end{exemple}

\subsection{Équivalence}

\begin{definition}
On écrit
\[
P\Longleftrightarrow Q
\]
si $P\Rightarrow Q$ et $Q\Rightarrow P$.
\end{definition}

\begin{methode}
Pour démontrer une équivalence $P\Longleftrightarrow Q$, on démontre généralement deux implications :
\[
P\Rightarrow Q,
\qquad
Q\Rightarrow P.
\]
\end{methode}

\subsection{Quantificateurs}

\begin{definition}
Les deux quantificateurs fondamentaux sont :
\[
\forall : \text{ pour tout},
\qquad
\exists : \text{ il existe}.
\]
\end{definition}

\begin{exemple}
\[
\forall x\in\R,\ x^2\geq 0
\]
signifie : pour tout réel $x$, on a $x^2\geq 0$.

\[
\exists x\in\R,\ x^2=2
\]
signifie : il existe au moins un réel $x$ tel que $x^2=2$.
\end{exemple}

\subsection{Ordre des quantificateurs}

\begin{remarque}
L'ordre des quantificateurs est essentiel.
\end{remarque}

\begin{exemple}
La proposition
\[
\forall x\in\R,\ \exists y\in\R,\ y>x
\]
est vraie : pour chaque réel $x$, on peut trouver un réel plus grand que lui.

Mais
\[
\exists y\in\R,\ \forall x\in\R,\ y>x
\]
est fausse : il n'existe pas de réel plus grand que tous les autres.
\end{exemple}

\subsection{Négation des quantificateurs}

\begin{proposition}
On a les règles fondamentales :
\[
\neg(\forall x\in E,\ P(x))
\Longleftrightarrow
\exists x\in E,\ \neg P(x),
\]
et
\[
\neg(\exists x\in E,\ P(x))
\Longleftrightarrow
\forall x\in E,\ \neg P(x).
\]
\end{proposition}

\begin{exemple}
La négation de
\[
\forall x\in\R,\ f(x)>0
\]
est
\[
\exists x\in\R,\ f(x)\leq 0.
\]
\end{exemple}

\begin{exemple}
La négation de
\[
\forall \eps>0,\ \exists \eta>0,\ \forall x\in E,\ d(x,a)<\eta \Rightarrow d(f(x),f(a))<\eps
\]
est
\[
\exists \eps>0,\ \forall \eta>0,\ \exists x\in E,\ d(x,a)<\eta
\quad \text{et}\quad
d(f(x),f(a))\geq \eps.
\]
Cette forme est typique lorsqu'on veut nier la continuité.
\end{exemple}

%================================================
\section{Ensembles et opérations ensemblistes}

\subsection{Ensembles, appartenance, inclusion}

\begin{definition}
Si $E$ est un ensemble et $x$ un objet, on écrit
\[
x\in E
\]
pour dire que $x$ appartient à $E$.
\end{definition}

\begin{definition}
Si $A$ et $B$ sont deux parties d'un ensemble $E$, on dit que $A$ est incluse dans $B$, et on écrit
\[
A\subset B,
\]
si
\[
\forall x\in E,\ x\in A \Rightarrow x\in B.
\]
\end{definition}

\begin{remarque}
Pour démontrer une inclusion $A\subset B$, on prend un élément quelconque $x\in A$ et on démontre que $x\in B$.
\end{remarque}

\subsection{Égalité d'ensembles}

\begin{proposition}
Pour montrer que deux ensembles $A$ et $B$ sont égaux, on montre souvent :
\[
A\subset B
\qquad \text{et} \qquad
B\subset A.
\]
\end{proposition}

\subsection{Réunion et intersection}

\begin{definition}
Soient $A,B\subset E$.
\[
A\cup B=\{x\in E,\ x\in A \text{ ou } x\in B\},
\]
\[
A\cap B=\{x\in E,\ x\in A \text{ et } x\in B\}.
\]
\end{definition}

\begin{definition}
Pour une famille de parties $(A_i)_{i\in I}$, on définit :
\[
\bigcup_{i\in I}A_i
=
\{x\in E,\ \exists i\in I,\ x\in A_i\},
\]
\[
\bigcap_{i\in I}A_i
=
\{x\in E,\ \forall i\in I,\ x\in A_i\}.
\]
\end{definition}

\subsection{Complémentaire}

\begin{definition}
Si $A\subset E$, le complémentaire de $A$ dans $E$ est
\[
A^c=E\setminus A=\{x\in E,\ x\notin A\}.
\]
\end{definition}

\subsection{Lois de De Morgan}

\begin{proposition}
Pour une famille $(A_i)_{i\in I}$ de parties de $E$, on a :
\[
\left(\bigcup_{i\in I} A_i\right)^c
=
\bigcap_{i\in I} A_i^c,
\]
et
\[
\left(\bigcap_{i\in I} A_i\right)^c
=
\bigcup_{i\in I} A_i^c.
\]
\end{proposition}

\begin{remarque}
Ces formules sont fondamentales en topologie, car les fermés sont définis comme les complémentaires des ouverts.
\end{remarque}

%================================================
\section{Applications, images et images réciproques}

\subsection{Applications}

\begin{definition}
Une application
\[
f:E\to F
\]
associe à tout élément $x\in E$ un unique élément $f(x)\in F$.
\end{definition}

\subsection{Image directe}

\begin{definition}
Si $A\subset E$, l'image directe de $A$ par $f$ est
\[
f(A)=\{f(x),\ x\in A\}.
\]
\end{definition}

\subsection{Image réciproque}

\begin{definition}
Si $B\subset F$, l'image réciproque de $B$ par $f$ est
\[
f^{-1}(B)=\{x\in E,\ f(x)\in B\}.
\]
\end{definition}

\begin{remarque}
Attention : $f^{-1}(B)$ est défini même si $f$ n'est pas bijective. Il ne s'agit pas forcément de l'application réciproque.
\end{remarque}

\subsection{Propriétés fondamentales}

\begin{proposition}
Pour toute famille $(B_i)_{i\in I}$ de parties de $F$, on a :
\[
f^{-1}\left(\bigcup_{i\in I}B_i\right)
=
\bigcup_{i\in I}f^{-1}(B_i),
\]
\[
f^{-1}\left(\bigcap_{i\in I}B_i\right)
=
\bigcap_{i\in I}f^{-1}(B_i),
\]
\[
f^{-1}(B^c)=\left(f^{-1}(B)\right)^c.
\]
\end{proposition}

\begin{remarque}
Ces propriétés expliquent pourquoi la continuité topologique se définit avec les images réciproques d'ouverts, et non avec les images directes.
\end{remarque}

\subsection{Injection, surjection, bijection}

\begin{definition}
Une application $f:E\to F$ est :
\begin{itemize}
    \item \textbf{injective} si
    \[
    \forall x,y\in E,\ f(x)=f(y)\Rightarrow x=y ;
    \]
    \item \textbf{surjective} si
    \[
    \forall z\in F,\ \exists x\in E,\ f(x)=z ;
    \]
    \item \textbf{bijective} si elle est injective et surjective.
\end{itemize}
\end{definition}

\begin{remarque}
Un homéomorphisme est une bijection continue dont la réciproque est continue.
C'est l'analogue topologique d'un isomorphisme.
\end{remarque}

%================================================
\section{Relations d'équivalence et ensembles quotients}

\subsection{Relations d'équivalence}

\begin{definition}
Une relation $\mathcal R$ sur un ensemble $E$ est une relation d'équivalence si elle est :
\begin{itemize}
    \item réflexive :
    \[
    \forall x\in E,\ x\mathcal R x ;
    \]
    \item symétrique :
    \[
    \forall x,y\in E,\ x\mathcal R y\Rightarrow y\mathcal R x ;
    \]
    \item transitive :
    \[
    \forall x,y,z\in E,\ x\mathcal R y \text{ et } y\mathcal R z
    \Rightarrow x\mathcal R z.
    \]
\end{itemize}
\end{definition}

\subsection{Classes d'équivalence}

\begin{definition}
La classe d'équivalence d'un élément $x\in E$ est
\[
[x]=\{y\in E,\ y\mathcal R x\}.
\]
\end{definition}

\begin{proposition}
Les classes d'équivalence forment une partition de $E$ :
\begin{itemize}
    \item elles sont non vides ;
    \item deux classes sont soit égales, soit disjointes ;
    \item leur réunion vaut $E$.
\end{itemize}
\end{proposition}

\subsection{Ensemble quotient}

\begin{definition}
L'ensemble quotient de $E$ par $\mathcal R$ est l'ensemble des classes d'équivalence :
\[
E/\mathcal R=\{[x],\ x\in E\}.
\]
\end{definition}

\begin{exemple}
Sur $[0,1]$, on peut identifier $0$ et $1$.
On obtient intuitivement un cercle :
\[
[0,1]/(0\sim 1)\simeq \mathbb S^1.
\]
\end{exemple}

\begin{remarque}
Les quotients sont essentiels en topologie : ils permettent de construire un cercle, un cylindre, un tore ou encore une bouteille de Klein en identifiant certains points.
\end{remarque}

%================================================
\section{Suites réelles et limites}

\subsection{Définition d'une suite}

\begin{definition}
Une suite réelle est une application
\[
u:\N\to\R.
\]
On note généralement $u_n$ au lieu de $u(n)$.
\end{definition}

\subsection{Convergence}

\begin{definition}
Une suite $(u_n)$ converge vers $\ell\in\R$ si
\[
\forall \eps>0,\ \exists N\in\N,\ \forall n\geq N,\ |u_n-\ell|<\eps.
\]
On note alors
\[
u_n\longrightarrow \ell.
\]
\end{definition}

\begin{remarque}
La topologie généralise précisément cette idée :
\[
u_n\to \ell
\]
signifie qu'à partir d'un certain rang, $u_n$ appartient à tout voisinage de $\ell$.
\end{remarque}

\subsection{Unicité de la limite}

\begin{proposition}
Si une suite réelle converge, alors sa limite est unique.
\end{proposition}

\begin{proof}
Supposons que $u_n\to \ell$ et $u_n\to m$ avec $\ell\neq m$.
Posons
\[
\eps=\frac{|\ell-m|}{3}>0.
\]
À partir d'un certain rang,
\[
|u_n-\ell|<\eps
\qquad\text{et}\qquad
|u_n-m|<\eps.
\]
Alors, par inégalité triangulaire,
\[
|\ell-m|
\leq |\ell-u_n|+|u_n-m|
<2\eps
=
\frac{2}{3}|\ell-m|,
\]
ce qui est impossible. Donc $\ell=m$.
\end{proof}

\subsection{Suites extraites}

\begin{definition}
Une suite extraite de $(u_n)$ est une suite de la forme
\[
u_{\varphi(n)},
\]
où $\varphi:\N\to\N$ est strictement croissante.
\end{definition}

\begin{proposition}
Si $u_n\to \ell$, alors toute suite extraite converge aussi vers $\ell$.
\end{proposition}

\subsection{Valeurs d'adhérence}

\begin{definition}
Un réel $\ell$ est une valeur d'adhérence de $(u_n)$ s'il existe une suite extraite $u_{\varphi(n)}$ telle que
\[
u_{\varphi(n)}\to \ell.
\]
\end{definition}

\begin{exemple}
La suite $u_n=(-1)^n$ admet deux valeurs d'adhérence :
\[
-1
\qquad\text{et}\qquad
1.
\]
\end{exemple}

%================================================
\section{Bornes, intervalles et propriétés de $\R$}

\subsection{Majorants, minorants}

\begin{definition}
Soit $A\subset\R$.
\begin{itemize}
    \item $M$ est un majorant de $A$ si
    \[
    \forall x\in A,\ x\leq M.
    \]
    \item $m$ est un minorant de $A$ si
    \[
    \forall x\in A,\ m\leq x.
    \]
\end{itemize}
\end{definition}

\subsection{Borne supérieure}

\begin{definition}
Si $A\subset\R$ est non vide et majorée, sa borne supérieure, notée $\sup A$, est le plus petit des majorants de $A$.
\end{definition}

\begin{proposition}
Si $s=\sup A$, alors :
\[
\forall x\in A,\ x\leq s,
\]
et
\[
\forall \eps>0,\ \exists x\in A,\ s-\eps<x\leq s.
\]
\end{proposition}

\subsection{Intervalles}

\begin{definition}
Une partie $I\subset\R$ est un intervalle si
\[
\forall a,b\in I,\ \forall x\in\R,\ a\leq x\leq b \Rightarrow x\in I.
\]
\end{definition}

\begin{remarque}
Les intervalles sont exactement les parties connexes de $\R$. Cette idée devient un théorème important en topologie.
\end{remarque}

%================================================
\section{Espaces vectoriels et normes}

\subsection{Espaces vectoriels}

\begin{definition}
Un espace vectoriel sur $\K=\R$ ou $\C$ est un ensemble $E$ muni :
\begin{itemize}
    \item d'une addition $(x,y)\mapsto x+y$ ;
    \item d'une multiplication externe $(\lambda,x)\mapsto \lambda x$ ;
\end{itemize}
vérifiant les règles usuelles de calcul vectoriel.
\end{definition}

\subsection{Norme}

\begin{definition}
Soit $E$ un espace vectoriel sur $\K$.
Une norme sur $E$ est une application
\[
\norme{\cdot}:E\to\R_+
\]
telle que, pour tous $x,y\in E$ et tout $\lambda\in\K$ :
\begin{itemize}
    \item $\norme{x}=0 \Longleftrightarrow x=0$ ;
    \item $\norme{\lambda x}=|\lambda|\norme{x}$ ;
    \item $\norme{x+y}\leq \norme{x}+\norme{y}$.
\end{itemize}
\end{definition}

\begin{exemple}
Sur $\R^n$, on définit :
\[
\norme{x}_1=\sum_{i=1}^n |x_i|,
\]
\[
\norme{x}_2=\left(\sum_{i=1}^n x_i^2\right)^{1/2},
\]
\[
\norme{x}_\infty=\max_{1\leq i\leq n}|x_i|.
\]
\end{exemple}

\subsection{Inégalité triangulaire inversée}

\begin{proposition}
Dans un espace vectoriel normé,
\[
\left|\norme{x}-\norme{y}\right|
\leq
\norme{x-y}.
\]
\end{proposition}

\begin{proof}
Par inégalité triangulaire,
\[
\norme{x}=\norme{x-y+y}\leq \norme{x-y}+\norme{y},
\]
donc
\[
\norme{x}-\norme{y}\leq \norme{x-y}.
\]
En échangeant $x$ et $y$, on obtient
\[
\norme{y}-\norme{x}\leq \norme{x-y}.
\]
D'où le résultat.
\end{proof}

%================================================
\section{Distances et espaces métriques}

\subsection{Distance}

\begin{definition}
Une distance sur un ensemble $E$ est une application
\[
d:E\times E\to\R_+
\]
telle que, pour tous $x,y,z\in E$ :
\begin{itemize}
    \item $d(x,y)=0\Longleftrightarrow x=y$ ;
    \item $d(x,y)=d(y,x)$ ;
    \item $d(x,z)\leq d(x,y)+d(y,z)$.
\end{itemize}
\end{definition}

\begin{definition}
Un couple $(E,d)$ est appelé espace métrique.
\end{definition}

\subsection{Distance associée à une norme}

\begin{proposition}
Si $(E,\norme{\cdot})$ est un espace vectoriel normé, alors
\[
d(x,y)=\norme{x-y}
\]
définit une distance sur $E$.
\end{proposition}

\subsection{Boules ouvertes et fermées}

\begin{definition}
Dans un espace métrique $(E,d)$, la boule ouverte de centre $a$ et de rayon $r>0$ est
\[
B(a,r)=\{x\in E,\ d(x,a)<r\}.
\]
La boule fermée de centre $a$ et de rayon $r\geq 0$ est
\[
\overline B(a,r)=\{x\in E,\ d(x,a)\leq r\}.
\]
\end{definition}

\begin{exemple}
Dans $\R$ muni de la distance usuelle,
\[
B(a,r)=]a-r,a+r[.
\]
\end{exemple}

\subsection{Suites dans un espace métrique}

\begin{definition}
Une suite $(x_n)$ d'un espace métrique $(E,d)$ converge vers $a\in E$ si
\[
\forall \eps>0,\ \exists N\in\N,\ \forall n\geq N,\ d(x_n,a)<\eps.
\]
\end{definition}

\begin{remarque}
Cette définition est exactement la définition usuelle dans $\R^n$, avec $d(x,a)=\norme{x-a}$.
\end{remarque}

%================================================
\section{Ouverts et fermés dans un espace métrique}

\subsection{Ouverts}

\begin{definition}
Une partie $U\subset E$ est ouverte si
\[
\forall x\in U,\ \exists r>0,\ B(x,r)\subset U.
\]
\end{definition}

\begin{exemple}
Dans $\R$, les intervalles ouverts $]a,b[$ sont des ouverts.
\end{exemple}

\begin{exemple}
Dans $\R^2$, un disque ouvert est un ouvert.
\end{exemple}

\subsection{Fermés}

\begin{definition}
Une partie $F\subset E$ est fermée si son complémentaire $E\setminus F$ est ouvert.
\end{definition}

\begin{proposition}
Dans un espace métrique, $F$ est fermé si et seulement si pour toute suite $(x_n)$ d'éléments de $F$ convergeant vers $x\in E$, on a $x\in F$.
\end{proposition}

\begin{proof}
On donne l'idée de la preuve.

Si $F$ est fermé et si $x\notin F$, alors $F^c$ est ouvert. Comme $x\in F^c$, il existe $r>0$ tel que
\[
B(x,r)\subset F^c.
\]
Une suite de points de $F$ ne peut donc pas converger vers $x$.

Réciproquement, si $F$ n'est pas fermé, alors $F^c$ n'est pas ouvert. Il existe donc un point $x\in F^c$ tel que toute boule centrée en $x$ rencontre $F$. On construit alors une suite $(x_n)$ dans $F$ vérifiant
\[
d(x_n,x)<\frac1n.
\]
Donc $x_n\to x$, mais $x\notin F$.
\end{proof}

\subsection{Adhérence}

\begin{definition}
L'adhérence d'une partie $A\subset E$, notée $\overline A$, est l'ensemble des points $x\in E$ tels que toute boule ouverte centrée en $x$ rencontre $A$ :
\[
x\in\overline A
\Longleftrightarrow
\forall r>0,\ B(x,r)\cap A\neq \varnothing.
\]
\end{definition}

\begin{proposition}
Dans un espace métrique,
\[
x\in\overline A
\Longleftrightarrow
\exists (a_n)\in A^\N,\ a_n\to x.
\]
\end{proposition}

\subsection{Intérieur}

\begin{definition}
L'intérieur de $A$, noté $\mathring A$, est l'ensemble des points $x\in A$ tels qu'il existe une boule centrée en $x$ incluse dans $A$ :
\[
x\in\mathring A
\Longleftrightarrow
\exists r>0,\ B(x,r)\subset A.
\]
\end{definition}

\subsection{Frontière}

\begin{definition}
La frontière de $A$ est
\[
\partial A=\overline A\setminus \mathring A.
\]
\end{definition}

\begin{remarque}
Un point est sur la frontière de $A$ si toute boule autour de ce point rencontre à la fois $A$ et son complémentaire.
\end{remarque}

%================================================
\section{Continuité}

\subsection{Continuité en un point}

\begin{definition}
Soient $(E,d_E)$ et $(F,d_F)$ deux espaces métriques.
Une application $f:E\to F$ est continue en $a\in E$ si
\[
\forall \eps>0,\ \exists \eta>0,\ \forall x\in E,\ d_E(x,a)<\eta
\Rightarrow
d_F(f(x),f(a))<\eps.
\]
\end{definition}

\subsection{Continuité globale}

\begin{definition}
$f$ est continue sur $E$ si elle est continue en tout point de $E$.
\end{definition}

\subsection{Caractérisation séquentielle}

\begin{proposition}
Dans les espaces métriques, $f:E\to F$ est continue en $a$ si et seulement si
\[
x_n\to a \Rightarrow f(x_n)\to f(a).
\]
\end{proposition}

\begin{remarque}
Cette caractérisation est très utile en prépa, mais en topologie générale elle ne suffit plus toujours. Certains espaces topologiques ne sont pas entièrement décrits par les suites.
\end{remarque}

\subsection{Continuité et ouverts}

\begin{theoreme}
Soient $(E,d_E)$ et $(F,d_F)$ deux espaces métriques.
Une application $f:E\to F$ est continue si et seulement si pour tout ouvert $V$ de $F$, $f^{-1}(V)$ est un ouvert de $E$.
\end{theoreme}

\begin{remarque}
Cette propriété devient la définition de la continuité en topologie générale.
\end{remarque}

\subsection{Continuité uniforme}

\begin{definition}
Une application $f:E\to F$ est uniformément continue si
\[
\forall \eps>0,\ \exists \eta>0,\ \forall x,y\in E,\ d_E(x,y)<\eta
\Rightarrow
d_F(f(x),f(y))<\eps.
\]
\end{definition}

\begin{remarque}
La différence avec la continuité simple est que $\eta$ ne dépend pas du point $a$, mais seulement de $\eps$.
\end{remarque}

\subsection{Applications lipschitziennes}

\begin{definition}
Une application $f:E\to F$ est lipschitzienne s'il existe $K\geq 0$ tel que
\[
\forall x,y\in E,\ d_F(f(x),f(y))\leq Kd_E(x,y).
\]
\end{definition}

\begin{proposition}
Toute application lipschitzienne est uniformément continue, donc continue.
\end{proposition}

%================================================
\section{Applications linéaires continues}

\subsection{Applications linéaires}

\begin{definition}
Soient $E$ et $F$ deux espaces vectoriels sur $\K$.
Une application $u:E\to F$ est linéaire si
\[
\forall x,y\in E,\ \forall \lambda,\mu\in\K,\ 
u(\lambda x+\mu y)=\lambda u(x)+\mu u(y).
\]
\end{definition}

\subsection{Continuité des applications linéaires}

\begin{theoreme}
Soient $(E,\norme{\cdot}_E)$ et $(F,\norme{\cdot}_F)$ deux espaces vectoriels normés.
Pour une application linéaire $u:E\to F$, les propriétés suivantes sont équivalentes :
\begin{enumerate}
    \item $u$ est continue sur $E$ ;
    \item $u$ est continue en $0$ ;
    \item il existe $C\geq 0$ tel que
    \[
    \forall x\in E,\ \norme{u(x)}_F\leq C\norme{x}_E ;
    \]
    \item $u$ est lipschitzienne.
\end{enumerate}
\end{theoreme}

\begin{proof}
On donne l'idée principale.

Si $u$ est continue en $0$, alors pour $\eps=1$, il existe $\eta>0$ tel que
\[
\norme{x}_E<\eta \Rightarrow \norme{u(x)}_F<1.
\]
Pour $x\neq 0$, on applique cette propriété à
\[
\frac{\eta}{2}\frac{x}{\norme{x}_E}.
\]
On obtient
\[
\norme{u(x)}_F
\leq
\frac{2}{\eta}\norme{x}_E.
\]
Donc $u$ est bornée par une constante fois $\norme{x}_E$, ce qui donne la continuité lipschitzienne.
\end{proof}

\subsection{Cas de la dimension finie}

\begin{theoreme}
Toute application linéaire entre espaces vectoriels normés de dimension finie est continue.
\end{theoreme}

\begin{remarque}
En dimension infinie, ce résultat est faux : il existe des applications linéaires non continues.
\end{remarque}

%================================================
\section{Équivalence des normes en dimension finie}

\begin{definition}
Deux normes $N_1$ et $N_2$ sur un espace vectoriel $E$ sont équivalentes s'il existe deux constantes $c,C>0$ telles que
\[
\forall x\in E,\ cN_1(x)\leq N_2(x)\leq C N_1(x).
\]
\end{definition}

\begin{theoreme}
Sur un espace vectoriel de dimension finie, toutes les normes sont équivalentes.
\end{theoreme}

\begin{remarque}
C'est un résultat très important : en dimension finie, la topologie ne dépend pas du choix de la norme.
\end{remarque}

\begin{exemple}
Sur $\R^n$, les normes $\norme{\cdot}_1$, $\norme{\cdot}_2$ et $\norme{\cdot}_\infty$ sont équivalentes.
On a par exemple
\[
\norme{x}_\infty\leq \norme{x}_2\leq \sqrt n \norme{x}_\infty.
\]
\end{exemple}

%================================================
\section{Compacité en analyse réelle}

\subsection{Définition intuitive}

\begin{remarque}
En prépa, on rencontre souvent la compacité sous la forme suivante :
dans $\R^n$, les compacts sont exactement les fermés bornés.
\end{remarque}

\begin{theoreme}[Heine-Borel]
Une partie $K\subset\R^n$ est compacte si et seulement si elle est fermée et bornée.
\end{theoreme}

\subsection{Caractérisation séquentielle}

\begin{theoreme}[Bolzano-Weierstrass]
Une partie $K\subset\R^n$ est compacte si et seulement si toute suite d'éléments de $K$ admet une suite extraite convergeant vers un élément de $K$.
\end{theoreme}

\subsection{Image continue d'un compact}

\begin{theoreme}
Si $K$ est compact et si $f:K\to F$ est continue, alors $f(K)$ est compact.
\end{theoreme}

\begin{corollaire}
Si $K$ est compact et si $f:K\to\R$ est continue, alors $f$ est bornée et atteint ses bornes :
\[
\exists x_{\min},x_{\max}\in K,\quad
f(x_{\min})=\min_K f,
\qquad
f(x_{\max})=\max_K f.
\]
\end{corollaire}

\begin{remarque}
En topologie générale, la compacité sera définie non pas avec les suites, mais avec les recouvrements par ouverts.
\end{remarque}

%================================================
\section{Connexité}

\subsection{Idée intuitive}

\begin{remarque}
Un espace est connexe s'il est \og en un seul morceau \fg.
Par exemple, un intervalle de $\R$ est connexe, mais
\[
[0,1]\cup [2,3]
\]
ne l'est pas.
\end{remarque}

\subsection{Connexité dans $\R$}

\begin{theoreme}
Les parties connexes de $\R$ sont exactement les intervalles.
\end{theoreme}

\subsection{Image continue d'un connexe}

\begin{theoreme}
L'image continue d'un connexe est connexe.
\end{theoreme}

\begin{corollaire}[Théorème des valeurs intermédiaires]
Si $I$ est un intervalle de $\R$ et si $f:I\to\R$ est continue, alors $f(I)$ est un intervalle.
\end{corollaire}

\subsection{Connexité par arcs}

\begin{definition}
Une partie $A$ d'un espace vectoriel normé est connexe par arcs si pour tous $x,y\in A$, il existe une application continue
\[
\gamma:[0,1]\to A
\]
telle que
\[
\gamma(0)=x,
\qquad
\gamma(1)=y.
\]
\end{definition}

\begin{proposition}
Toute partie connexe par arcs est connexe.
\end{proposition}

\begin{exemple}
Tout convexe d'un espace vectoriel normé est connexe par arcs, donc connexe.
\end{exemple}

\begin{proof}
Si $C$ est convexe et $x,y\in C$, le chemin
\[
\gamma(t)=(1-t)x+ty
\]
reste dans $C$ pour tout $t\in[0,1]$.
\end{proof}

%================================================
\section{Suites de Cauchy et complétude}

\subsection{Suites de Cauchy}

\begin{definition}
Soit $(E,d)$ un espace métrique.
Une suite $(x_n)$ est de Cauchy si
\[
\forall \eps>0,\ \exists N\in\N,\ \forall p,q\geq N,\ d(x_p,x_q)<\eps.
\]
\end{definition}

\begin{remarque}
Une suite convergente est toujours de Cauchy.
\end{remarque}

\begin{proposition}
Toute suite convergente dans un espace métrique est de Cauchy.
\end{proposition}

\begin{proof}
Supposons $x_n\to \ell$.
Soit $\eps>0$.
Il existe $N$ tel que, pour tout $n\geq N$,
\[
d(x_n,\ell)<\frac{\eps}{2}.
\]
Alors, pour $p,q\geq N$,
\[
d(x_p,x_q)
\leq d(x_p,\ell)+d(\ell,x_q)
<\eps.
\]
\end{proof}

\subsection{Espaces complets}

\begin{definition}
Un espace métrique $(E,d)$ est complet si toute suite de Cauchy d'éléments de $E$ converge vers un élément de $E$.
\end{definition}

\begin{exemple}
$\R$ est complet.
\end{exemple}

\begin{exemple}
$\Q$ muni de la distance usuelle n'est pas complet.
Par exemple, une suite rationnelle convergeant vers $\sqrt 2$ est de Cauchy dans $\Q$, mais ne converge pas dans $\Q$.
\end{exemple}

\subsection{Complétude et fermés}

\begin{proposition}
Si $(E,d)$ est complet et si $F\subset E$ est fermé, alors $F$ est complet.
\end{proposition}

\begin{proof}
Soit $(x_n)$ une suite de Cauchy dans $F$.
Comme $E$ est complet, $(x_n)$ converge vers un certain $x\in E$.
Comme $F$ est fermé et que tous les $x_n$ sont dans $F$, on a $x\in F$.
Donc $F$ est complet.
\end{proof}

%================================================
\section{Topologie : passage de l'analyse à l'abstraction}

\subsection{Pourquoi généraliser ?}

En analyse de prépa, on travaille surtout dans :
\[
\R,\quad \R^n,\quad \C,\quad \text{ou des espaces vectoriels normés}.
\]
Dans ces espaces, la notion de proximité est donnée par une distance ou une norme.

Mais en topologie générale, on veut parler de continuité, de voisinage, de convergence, de compacité et de connexité sans forcément disposer d'une distance.

\subsection{Idée centrale}

\begin{definition}
Une topologie sur un ensemble $E$ est une famille $\tau\subset \mathcal P(E)$ telle que :
\begin{itemize}
    \item $\varnothing\in\tau$ et $E\in\tau$ ;
    \item toute réunion d'éléments de $\tau$ appartient à $\tau$ ;
    \item toute intersection finie d'éléments de $\tau$ appartient à $\tau$.
\end{itemize}
Les éléments de $\tau$ sont appelés ouverts.
\end{definition}

\begin{remarque}
Dans un espace métrique, les ouverts définis par les boules vérifient exactement ces trois propriétés.
\end{remarque}

\subsection{Topologie associée à une distance}

\begin{definition}
Si $(E,d)$ est un espace métrique, on définit une topologie $\tau_d$ par :
\[
U\in\tau_d
\Longleftrightarrow
\forall x\in U,\ \exists r>0,\ B(x,r)\subset U.
\]
\end{definition}

\begin{remarque}
Ainsi, les espaces métriques sont des cas particuliers d'espaces topologiques.
\end{remarque}

%================================================
\section{Voisinages}

\begin{definition}
Soit $(E,\tau)$ un espace topologique et $x\in E$.
Une partie $V\subset E$ est un voisinage de $x$ s'il existe un ouvert $U\in\tau$ tel que
\[
x\in U\subset V.
\]
\end{definition}

\begin{remarque}
Dans un espace métrique, dire que $V$ est un voisinage de $x$ signifie que $V$ contient une boule ouverte centrée en $x$.
\end{remarque}

\begin{exemple}
Dans $\R$, $[0,2]$ est un voisinage de $1$, car il contient par exemple l'ouvert
\[
]1/2,3/2[.
\]
Mais $[0,2]$ n'est pas un voisinage de $0$ dans $\R$, car il ne contient aucun intervalle ouvert centré en $0$.
\end{exemple}

%================================================
\section{Bases d'ouverts}

\begin{definition}
Une famille $\mathcal B$ d'ouverts d'un espace topologique $(E,\tau)$ est une base d'ouverts si tout ouvert de $E$ est réunion d'éléments de $\mathcal B$.
\end{definition}

\begin{exemple}
Dans un espace métrique, les boules ouvertes forment une base d'ouverts.
\end{exemple}

\begin{exemple}
Dans $\R$, les intervalles ouverts forment une base d'ouverts.
\end{exemple}

\begin{remarque}
Les bases d'ouverts permettent de décrire une topologie sans avoir besoin de lister tous les ouverts.
\end{remarque}

%================================================
\section{Séparation}

\begin{definition}
Un espace topologique $(E,\tau)$ est séparé, ou de Hausdorff, si pour tous $x,y\in E$ avec $x\neq y$, il existe deux ouverts $U,V$ tels que
\[
x\in U,\qquad y\in V,
\qquad U\cap V=\varnothing.
\]
\end{definition}

\begin{proposition}
Tout espace métrique est séparé.
\end{proposition}

\begin{proof}
Soient $x\neq y$.
Alors $d(x,y)>0$.
Posons
\[
r=\frac{d(x,y)}{3}.
\]
Alors les boules
\[
B(x,r)
\qquad\text{et}\qquad
B(y,r)
\]
sont disjointes.
\end{proof}

\begin{remarque}
Dans un espace séparé, une suite ne peut pas avoir deux limites différentes.
\end{remarque}

%================================================
\section{Topologie induite, produit et quotient}

\subsection{Topologie induite}

\begin{definition}
Soit $(E,\tau)$ un espace topologique et $A\subset E$.
La topologie induite sur $A$ est
\[
\tau_A=\{U\cap A,\ U\in\tau\}.
\]
\end{definition}

\begin{exemple}
Dans $\R$, l'intervalle $[0,1]$ hérite d'une topologie.
Par exemple,
\[
[0,1/2[
=
]-1,1/2[\cap [0,1]
\]
est ouvert dans $[0,1]$ pour la topologie induite, même s'il n'est pas ouvert dans $\R$.
\end{exemple}

\subsection{Topologie produit}

\begin{definition}
Si $(E,\tau_E)$ et $(F,\tau_F)$ sont deux espaces topologiques, la topologie produit sur $E\times F$ est celle dont les ouverts de base sont les produits
\[
U\times V,
\qquad
U\in\tau_E,
\quad
V\in\tau_F.
\]
\end{definition}

\begin{exemple}
La topologie usuelle de $\R^2$ est la topologie produit de $\R\times\R$.
\end{exemple}

\subsection{Topologie quotient}

\begin{definition}
Soit $(E,\tau)$ un espace topologique et $\mathcal R$ une relation d'équivalence sur $E$.
On note
\[
\pi:E\to E/\mathcal R
\]
la projection canonique.

La topologie quotient sur $E/\mathcal R$ est définie par :
\[
U\subset E/\mathcal R \text{ est ouvert}
\Longleftrightarrow
\pi^{-1}(U)\text{ est ouvert dans }E.
\]
\end{definition}

\begin{remarque}
La topologie quotient est la bonne topologie pour rendre la projection canonique continue.
\end{remarque}

%================================================
\section{Résumé des prérequis à maîtriser}

Avant d'attaquer un cours de topologie générale, il faut être à l'aise avec :

\begin{enumerate}
    \item les quantificateurs $\forall$, $\exists$ ;
    \item les négations d'énoncés complexes ;
    \item les ensembles, réunions, intersections, complémentaires ;
    \item les images réciproques d'applications ;
    \item les relations d'équivalence et les quotients ;
    \item la convergence des suites ;
    \item les suites extraites et valeurs d'adhérence ;
    \item les espaces vectoriels normés ;
    \item les normes usuelles de $\R^n$ ;
    \item les distances et boules ouvertes ;
    \item les ouverts, fermés, adhérences, intérieurs, frontières ;
    \item la continuité métrique ;
    \item la continuité via image réciproque d'ouverts ;
    \item les applications linéaires continues ;
    \item l'équivalence des normes en dimension finie ;
    \item la compacité de $\R^n$ ;
    \item la connexité des intervalles ;
    \item la connexité par arcs ;
    \item les suites de Cauchy ;
    \item la complétude.
\end{enumerate}

%================================================
\section{Mini-fiche : dictionnaire prépa vers topologie générale}

\[
\begin{array}{|c|c|}
\hline
\textbf{Analyse de prépa} & \textbf{Topologie générale}\\
\hline
\text{Intervalle ouvert, boule ouverte} & \text{Ouvert abstrait}\\
\hline
\text{Boule autour de }x & \text{Voisinage de }x\\
\hline
|x_n-\ell|\to 0 & x_n\to \ell \text{ au sens des voisinages}\\
\hline
\text{Fonction continue avec }\eps,\eta
&
\text{Image réciproque d'un ouvert ouverte}\\
\hline
\text{Fermé stable par limite de suites}
&
\text{Complémentaire d'un ouvert}\\
\hline
\text{Compact = fermé borné dans }\R^n
&
\text{Recouvrement ouvert extractible fini}\\
\hline
\text{Intervalle}
&
\text{Connexe}\\
\hline
\text{Suite de Cauchy}
&
\text{Complétude métrique}\\
\hline
\text{Identifier }0\text{ et }1
&
\text{Topologie quotient}\\
\hline
\end{array}
\]

%================================================
\section{Exercices de vérification}

\subsection*{Exercice 1}
Écrire la négation de :
\[
\forall \eps>0,\ \exists N\in\N,\ \forall n\geq N,\ |u_n-\ell|<\eps.
\]

\subsection*{Exercice 2}
Montrer que si $A\subset B$, alors
\[
\overline A\subset \overline B
\]
dans un espace métrique.

\subsection*{Exercice 3}
Montrer que dans un espace métrique, toute boule ouverte est un ouvert.

\subsection*{Exercice 4}
Soit $f:E\to F$ une application continue entre espaces métriques.
Montrer que pour tout fermé $C$ de $F$, $f^{-1}(C)$ est fermé dans $E$.

\subsection*{Exercice 5}
Montrer que toute application linéaire de $\R^n$ dans $\R^p$ est continue.

\subsection*{Exercice 6}
Montrer que $[0,1]$ est compact dans $\R$.

\subsection*{Exercice 7}
Montrer que $]0,1[$ n'est pas compact.

\subsection*{Exercice 8}
Montrer que $[0,1]\cup[2,3]$ n'est pas connexe.

\subsection*{Exercice 9}
Montrer qu'un convexe de $\R^n$ est connexe par arcs.

\subsection*{Exercice 10}
Montrer que $\Q$ n'est pas complet pour la distance usuelle.

%================================================
\section{Corrigés rapides}

\subsection*{Correction 1}
La négation est :
\[
\exists \eps>0,\ \forall N\in\N,\ \exists n\geq N,\ |u_n-\ell|\geq \eps.
\]

\subsection*{Correction 2}
Soit $x\in\overline A$.
Alors toute boule centrée en $x$ rencontre $A$.
Comme $A\subset B$, toute boule centrée en $x$ rencontre aussi $B$.
Donc $x\in\overline B$.

\subsection*{Correction 3}
Soit $B(a,r)$ une boule ouverte.
Soit $x\in B(a,r)$.
Alors
\[
d(a,x)<r.
\]
Posons
\[
\rho=r-d(a,x)>0.
\]
Si $y\in B(x,\rho)$, alors
\[
d(a,y)\leq d(a,x)+d(x,y)<d(a,x)+\rho=r.
\]
Donc
\[
B(x,\rho)\subset B(a,r).
\]
Ainsi $B(a,r)$ est ouverte.

\subsection*{Correction 4}
Si $C$ est fermé, alors $F\setminus C$ est ouvert.
Comme $f$ est continue,
\[
f^{-1}(F\setminus C)
\]
est ouvert.
Mais
\[
f^{-1}(F\setminus C)=E\setminus f^{-1}(C).
\]
Donc $f^{-1}(C)$ est fermé.

\subsection*{Correction 5}
En dimension finie, toutes les normes sont équivalentes.
Si $u:\R^n\to\R^p$ est linéaire, elle est représentée par une matrice.
On peut majorer chaque coordonnée de $u(x)$ par une constante fois $\norme{x}$.
Il existe donc $C>0$ tel que
\[
\norme{u(x)}\leq C\norme{x}.
\]
Donc $u$ est lipschitzienne, donc continue.

\subsection*{Correction 6}
Dans $\R$, le segment $[0,1]$ est fermé et borné.
Par le théorème de Heine-Borel, il est compact.

\subsection*{Correction 7}
La suite
\[
u_n=\frac1n
\]
appartient à $]0,1[$ et converge vers $0$.
Or $0\notin ]0,1[$.
Donc $]0,1[$ n'est pas compact, car une suite d'éléments de cet intervalle n'admet pas nécessairement de sous-suite convergeant dans l'intervalle.

\subsection*{Correction 8}
On écrit
\[
[0,1]\cup[2,3]
=
\left([0,1]\cup[2,3]\right)\cap ]-\infty,3/2[
\ \cup\
\left([0,1]\cup[2,3]\right)\cap ]3/2,+\infty[.
\]
Ces deux ensembles sont ouverts dans la topologie induite, non vides et disjoints.
Donc l'ensemble n'est pas connexe.

\subsection*{Correction 9}
Soit $C$ un convexe de $\R^n$.
Pour $x,y\in C$, on pose
\[
\gamma(t)=(1-t)x+ty.
\]
Comme $C$ est convexe, $\gamma(t)\in C$ pour tout $t\in[0,1]$.
De plus, $\gamma$ est continue.
Donc $C$ est connexe par arcs.

\subsection*{Correction 10}
On peut construire une suite de rationnels $(q_n)$ convergeant vers $\sqrt2$ dans $\R$.
Alors $(q_n)$ est de Cauchy dans $\Q$.
Mais elle ne converge vers aucun rationnel, car sa limite dans $\R$ est $\sqrt2\notin\Q$.
Donc $\Q$ n'est pas complet.

\end{document}